\documentclass[10pt]{article}

\usepackage[utf8]{inputenc}

\usepackage[T1]{fontenc}

\usepackage{amsmath}

\usepackage{amsfonts}

\usepackage{amssymb}

\usepackage[version=4]{mhchem}

\usepackage{stmaryrd}

\usepackage{graphicx}

\usepackage[export]{adjustbox}

\graphicspath{ {./images/} }



\begin{document}

ность и существование рептения этой задачи как для уравнения (5), так и цля более общих уравнений доказаны при нек-рых условиях геометрич. характера на границу области $\Omega$, особенно на кривую $\sigma$. Общую смешаннуго задачу Бицадзе можно считать полностью исследованной в частном случае, когда кривая $\Gamma_{1}$ совпадает $\left(x_{0}=1\right)$ с выходящей из точки $B$ характеристикоІ̆ $B C$. Важным следствием корректности общей смешанной задати Бицадзе, напр. для уравнения (5), является тот факт, что для смешанных областей вида $\Omega$ Дирихле задача некорректна независимо от величины и формы области гиперболичности $\Omega$-.



Для доволњно широкого класса линейных уравнений



\[

k(y) u_{x x}+u_{y y} \frac{1}{1} a u_{x}+b u_{y}+c u=f

\]



установлено, что на корректность задачи Дирихле в соответствующих смешанных областях вида $\Omega$ существенное влияние может оказать коәффициент $a(x, y)$.



Смешанной задачей нового типа является з а д а ч а Ф р а н л л . Пусть односвязная область $\Omega$ ограничена: отрезком $A^{\prime} A,-1 \leqslant y \leqslant 1$, прямой $x=0$; гладкой кривоїі $\sigma$ с концами в точках $A(0,1)$ и $B(a, 0)$, расположенноїі в квадранте $x>0, y>0$; отрезком $C B: a_{1} \leqslant$ $\leqslant x \leqslant a$ прямой $y=0$ и проходящей через точки $A^{\prime}(0$, -1), $C\left(a_{1}, 0\right)$ характеристикой рассматриваемого С.т.у., напр. уравнения (4). Задача Франкля состоит в отыскании решения $u(x, y) \mathrm{C}$. т. у. в области $\Omega$, когда на $\sigma \bigcup C B$ задаются значения $u(x, y)$, а на $A^{\prime} A-$ условие



\[

\frac{\partial u}{\partial x}=0, u(0, y)-u(0,-y)=f(y),-1 \leqslant y \leqslant 1, x=0 .

\]



Эта задача исследована в основном для модельных С. т. У. Задача Франкля весьма полно решена для уравнения (5), если кривая $\sigma: x=x(s), y=y(s)$ такова, что $d y / d s \geqslant 0$, где $s$ - длина $\sigma$, отсчитываемая от точки $B(a, 0)$.



Сформулированные основные краевые задачи для С. т. у. 1-го рода с соответствующими изменениями перенесены на С. т. у. 2-го рода. Эти изменения вызваны тем, что задача Дирихле для эллиптич. уравнений с характеристич. вырождением на части границы не всегда является корректно поставленной.



В постановке краевых задач для уравнения (1) в смешанных областях вносится новый аспект, если линия $\delta$ изменения типа одновременно представляет собой линию вырождения порядка уравнения, напр. в случае уравнения



\[

y^{2} p u_{x x}+y u_{y y}+\beta u_{y}=0,

\]



где $p-$ натуральное число, а $\beta-$ постоянная, удовлетворяющая неравенству $1-2 p \leqslant 2 \beta<1$.



Для уравнений $(5),(6),(7)$, помимо указанных задач, исследован также ряд принципиально новых краевых задач, к-рые в основном характеризуются тем, что вся граница $\sigma \bigcup A C \bigcup B C$ области $\Omega$ (где ставится задача Трикоми) является носителем краевых условий: на кривой $\sigma$ задано, напр., условие Дирихле, а на $A C \bigcup B C-$ нек-рое нелокальное условие, поточечно связывающее значения искомого решения или его (дробной) производной определенного порядка. В частности, эти задачи включают простой пример корректной самосопряженной смешанной краевой задачи.



Изучаются краевые задачи для С. т. у. и систем в областях, содержапци внутри себя части нескольких линий вырождения типа или одну замкнутую параболич. линию.



Аналоги задачи Трикоми рассматривались и для нек-рого класса С. т. У. и систем с двумя независимыми переменными и уравнений высокого порядка. Значительные затруднения возникают при отыскании правильно поставленных задач для С. т. у. с многими независимыми иеременными. Тем не менее в әтом направлении получен ряд важных результатов. Установлено, что для уравнения



\[

\operatorname{sign} z \cdot u_{x x}+u_{y y}+u_{z z}=f(x, y, z),

\]



к-рое представляет собой простую модель С. т. у. с временны́ образом ориентированной плоскостью $z=0$ вырождения типа, корректно поставленной является следующая задача. Пусть $\Omega$ - конечная односвязная трехмерная область, ограниченная нек-рой кусочно гладкой поверхностью $z=f(x, y) \geqslant 0$ и характеристиками $S_{1}: x+x_{0}=\sqrt{y^{2}+z^{2}}, \quad S_{2}: x-x_{0}=\sqrt{y^{2}+z^{2}}$ уравнения (8). Требуется найти непрерывную в $\Omega$ функцию $u(x, y, z)$ с вепрерывными в $\Omega$ производными 1-го порядка, удовлетворяющую уравнению (8) в области $\Omega$ при $z \neq 0$ и обращающуюся в нуль на $\sigma$ и на одной из характеристик $S_{1}$ или $S_{2}$. Доказаны существование слабого и единственность сильного решения этої задачи и для более общего уравнения



\[

\operatorname{sign} x_{n} \cdot u_{x_{0} x_{0}}+\Delta_{x} u=f\left(x_{0}, x\right), x=\left(x_{1}, \ldots, x_{n}\right),

\]



где $\Delta_{x}$ - оператор Лапласа по переменным $x_{1}, \ldots, x_{n}$.



Для уравнения



\[

x_{0}^{2 m} \Delta_{x} u-x_{0} u_{x_{0} x_{0}}+(m-1 / 2) u_{x_{0}}=0

\]



с пространственно ориентированной гиперплоскостью $x_{0}=0$ вырождения тиша и порядка в смешанной области $\Omega$ специального вида поставлены и исследованы краевые задачи, в к-рых часть границы $\partial \Omega$, лежащая в полупространстве $x_{0}<0$, является носителем данных $u\left(x_{0}, x\right)$, а лежащая в полупространстве $x_{0}>0$ часть границы $\partial \Omega$ (характеристич. коноид уравнения (9)) носителем нек-рых интегральных средних для $u\left(x_{0}, x\right)$.



Исследовались и другие модельные С. т. у. в трехмерных ограниченных и неограниченных областях, в т. ч. уравнения



\[

z^{2 m+1} u_{x x}+u_{y y}+u_{z z}=0, z^{2 m+1}\left(u_{x x}+u_{y y}\right)+u_{z z}=0 .

\]



Найден также критерий единственности решения задачи Дирихле для широкого класса самосопряженных С. т. У. в цилиндрич. областях.



Лит.: [1] Б е р с Л., Математические вопросы дозвуковой и околозвуковой газовой динамики, пер. с англ., М., 1961; [2] Б И ц а д з е А. В., Уравнения смешанного типа, М., 1959; которых вырождается вдоль линии изменения типа, в сб.: Механика сплошной среды и родственные проблемы анализа, М. CССР", 1972, т. 205, № 1; [5] В е к у а И. Н., Обоб̆щенные аналитические функции, М., 1959; [6] К а р а т о п р а к л и е в



\begin{center}

\includegraphics[max width=\textwidth]{2023_07_09_d33f3059e745d3ab6217g-1(1)}

\end{center}



\includegraphics[max width=\textwidth, center]{2023_07_09_d33f3059e745d3ab6217g-1}

CCP. Серия физ.-мат. наук», 1969, № 1, с. 27-33; [9] С м и рн о в М. М., Уравнения смешанного типа, М., 1970; [10] С о Ј№2, с. $325-32$; [11] Т р и к о м и Ф., О линейных уравнениях в частных производных второго порядка смешанного типа, пер. с итал., М. - Л., 1947; [12] Ф р а н к л ь Ф. И., Избранные труды по газовой динамике, M., 1973; [13]F r i e d r i c h s K. О., "Communication Pure and Appl. Math.", 1958, v. 11, № 3, p., 1936, bd 26 A, № 3, s. $1-32$; [15] G e r m a i n P., B a de r R., СМЕІІАННОЕ ПРОИЗВЕДЕНИЕ $(a, b, c)$ в е к т ор о в $\boldsymbol{a}, \boldsymbol{b}, \boldsymbol{c}$ - скалярное произведение вектора $\boldsymbol{a}$ на векторное произведение векторов $\boldsymbol{b}$ и $c$ :



\[

(\boldsymbol{a}, \boldsymbol{b}, \boldsymbol{c})=(\boldsymbol{a},[\boldsymbol{b}, \boldsymbol{c}])

\]



См. Векторная алгебра.



СМЕІІАННЫЙ ПРОІЕСС АВТОРЕГРЕССИИСКОЛЬЗЯЩЕГО СРЕДНЕГО, А Р С С - П $\mathrm{p}$ о Ц е с с стационарный в широком смысле случайный процесс





\end{document}